\begin{abstract}
Automated program verification has seen significant advancements. 
These verifiers are typically constructed as a pipeline of translations: from high-level languages 
annotated with pre- and post-conditions, to intermediate verification languages (IVLs), 
and ultimately to verification conditions (VCs) encoded as Satisfiability Modulo Theories (SMT) formulas. 
Boogie is a prominent IVL whose verifier generates such VCs. However, the reliability of these 
verifiers depends entirely on their correct implementation, for which formal guarantees are rarely provided.

Prior work tackled this issue by generating certificates,
in the form of machine-checked proofs in Isabelle/HOL,
for individual verification runs.
This foundational framework supported a core 
subset of the Boogie language by instrumenting the verifier's codebase. Because the official Boogie 
repository has since undergone significant architectural updates, our first contribution modernizes 
the certificate generation module to ensure compatibility with the latest Boogie codebase.

Following a usage analysis, we identified Boogie maps as the most critical unsupported feature 
preventing the certification of complex programs from front-ends like Dafny. The primary contribution 
of this thesis is the design and extension of the formal foundational semantics to support a 
practical subset of Boogie maps. We extended the certificate generator accordingly and successfully 
certified verification runs for programs using maps, demonstrating the soundness of our approach. 
Finally, we discuss the advantages and limitations of our chosen semantic embedding and outline 
alternative approaches and directions for future work.
\end{abstract}